\chapter{Server interface}
\label{ch:server-interface}

The client reaches the database in two ways: by calling database functions and by reading the few tables
it is allowed to read. \Cref{fig:call-anatomy} follows one function call from the store to the tables and
back. This chapter is the reference for the interface; \cref{app:db-functions,app:tables} give the exact
signatures and table definitions.

\begin{figure}[ht]
  \centering
  % Anatomy of a database function call, with send_request as the example (Figure: call anatomy).
\begin{tikzpicture}[
    stage/.style={box, font=\sffamily\scriptsize, text width=25mm, minimum height=19mm, align=center},
    lab/.style={font=\sffamily\scriptsize, text=muted, inner sep=1.2pt, align=center},
  ]
  \node[stage, clientbox] (store) at (1.45,0) {\textbf{Store action}\\[2pt]\texttt{sayHi(person)}};
  \node[stage, accentbox] (api) at (4.65,0) {\textbf{API layer}\\[2pt]\texttt{supabase.rpc(}\\\texttt{'send\_request',}\\\texttt{\{target\})}};
  \node[stage, serverbox] (rest) at (7.85,0) {\textbf{Data API}\\[2pt]verifies the token;\\runs as role\\\texttt{authenticated}};
  \node[stage, serverbox, text width=27mm] (fn) at (11.15,0) {\textbf{\texttt{public.}}\\\textbf{\texttt{send\_request}}\\[2pt]\texttt{SECURITY DEFINER}\\\texttt{search\_path = ''}\\caller = \texttt{auth.uid()}};
  \node[stage, serverbox] (tables) at (14.45,0) {\textbf{Tables}\\[2pt]read and written as the owner, past row-level security};

  \draw[flow] (store) -- (api);
  \draw[flow] (api) -- (rest);
  \draw[flow] (rest) -- (fn);
  \draw[flow] (fn) -- (tables);

  \node[box, font=\sffamily\scriptsize, text width=40mm, align=flush left, anchor=north] (checks) at (12.65,-1.45)
    {\textbf{Checks, in order}\\[2pt]
     1. signed in; not the caller\\
     2. no block either way; target active\\
     3. caller not already chatting\\
     4. both live, target open, near\\
     5. under 5 requests in 10 minutes\\
     6. no other request waiting\\[2pt]
     then insert the \texttt{pending} row};
  \draw[flow] (fn.south) |- ($(checks.north)+(0,0.25)$) -- (checks.north);

  \node[stage, minimum height=12mm, text width=27mm, accentbox] (ok) at (7.85,-2.6) {success: JSON\\\texttt{\{id, expiresAt\}}};
  \node[stage, minimum height=12mm, text width=27mm, draw=alert, fill=alertsoft] (no) at (7.85,-4.3) {refusal: exception\\\texttt{'not nearby'}};
  \draw[flow] (checks.west |- ok) -- (ok);
  \draw[flow] (checks.west |- no) -- (no);
  \node[stage, minimum height=12mm, accentbox] (convert) at (4.65,-2.6) {typed result};
  \node[stage, minimum height=12mm, accentbox] (sentence) at (4.65,-4.3) {sentence from\\\texttt{SERVER\_REFUSALS}};
  \draw[flow] (ok) -- (convert);
  \draw[flow] (no) -- (sentence);
  \node[stage, minimum height=12mm, clientbox] (stateset) at (1.45,-2.6) {state updated;\\waiting screen};
  \node[stage, minimum height=12mm, clientbox] (alert) at (1.45,-4.3) {alert: ``They're not\\nearby any more.''};
  \draw[flow] (convert) -- (stateset);
  \draw[flow] (sentence) -- (alert);
\end{tikzpicture}

  \caption{Anatomy of a database function call, with \code{send_request} as the example. The function
    runs as its owner, takes the caller's identity from the token, and either writes after its checks or
    raises a short refusal that the API layer turns into a sentence.}
  \label{fig:call-anatomy}
\end{figure}

\section{Conventions}
\label{sec:function-conventions}

All database functions follow the same conventions, set out at the top of the migration that defines
them.

\begin{itemize}
  \item Functions in \code{public} are the API. They are declared \code{SECURITY DEFINER}, so they run
    as their owner and are not limited by row-level security, and \code{SET search_path = ''}, so every
    name inside them is schema-qualified and cannot be redirected by the caller.
  \item The caller is always \code{auth.uid()}, read from the session token. No function takes a user ID
    that it trusts as the caller's identity.
  \item Execution is revoked from the roles \code{public} and \code{anon}, so only a signed-in user
    (role \code{authenticated}) can call them. Helpers in \code{private} cannot be called through the
    Data API at all, because that schema is not exposed.
  \item A function refuses a request by raising an exception with a short, fixed message, for example
    \code{not nearby} or \code{too many requests}. The client translates each message into a sentence
    (\cref{app:refusals}).
  \item Functions that return a full profile return person cards (\cref{sec:person-card});
    \code{blocked_list} and an event card's host return a reduced object with no introduction, interests
    or note. Functions that return events return event cards. Both are JSON objects with camel-case keys
    that match the client's types, and the instants in them are milliseconds since the Unix epoch. The
    two message functions are the exception: they return the stored \code{messages} row, whose
    \texttt{created\_at} is an ISO~8601 string.
\end{itemize}

\section{Database functions}
\label{sec:functions}

\Cref{tab:functions} lists every function the client calls, in the order of the user's journey. The
\emph{Checks} column lists what the function verifies before it changes anything, with the refusal
message in parentheses where one applies.

\begingroup\tablefont
\begin{longtable}{@{}L{38mm}L{\dimexpr\linewidth-38mm-2\tabcolsep\relax}@{}}
  \caption{Database functions called by the client.}
  \label{tab:functions}\\
  \toprule
  \textsf{\textbf{Function}} & \textsf{\textbf{Effect and checks}} \\
  \midrule
  \endfirsthead
  \toprule
  \textsf{\textbf{Function}} & \textsf{\textbf{Effect and checks}} \\
  \midrule
  \endhead
  \bottomrule
  \endfoot
  \multicolumn{2}{@{}l}{\textsf{\textbf{Going live}}} \\[2pt]
  \code{go_green}\newline{\scriptsize lat, lng, accuracy, minutes, note}\newline{\scriptsize returns start and end times}
  & Starts or updates the caller's live session and returns its server times.
    Checks: signed in (\code{not signed in}); profile finished and adult (\code{profile not ready}) and
    active (\code{account restricted}); duration is 30, 60 or 120 (\code{invalid duration}); note of
    at most 60 characters (\code{note too long}); latitude and longitude in range (\code{invalid location}).
    Applies the relocation rules of \cref{sec:relocation-limits} under a per-user lock: a move of 150\,m
    or less keeps the stored position and start time; a larger move while live must not exceed 150\,m plus
    10\,km/h over the time since the position was set (\code{moving too fast}); going live when not live
    must be within 5\,km plus 1000\,km/h of the last known position (\code{location jumped}); and at most
    three location changes are allowed in ten minutes, shared with \code{set_map_position} (\code{too many
    location changes}). Stores the centre of the position's cell, the accuracy clamped to 0--50\,m, and the
    session as \code{chatting} if the caller has an accepted chat, otherwise \code{open}, ending
    \code{minutes} after now. The stored position also becomes the caller's last known position. \\
  \code{update_note}\newline{\scriptsize note}
  & Replaces the live note. Checks: an active profile, as for \code{go_green}; at most 60 characters
    (\code{note too long}). Does nothing if the caller is not live. \\
  \code{go_offline}
  & Withdraws the caller's waiting request (status \code{cancelled}), then deletes the caller's live
    session. Always succeeds. \\
  \code{nearby}\newline{\scriptsize returns person cards}
  & Returns up to 50 person cards of users who are live, \code{open}, not expired and active, in the
    caller's cell or a touching one (\cref{sec:location-privacy}), excluding the caller and anyone blocked
    in either direction. Returns nothing if the caller is not live. Never returns a position or a
    distance. At most 12 calls a minute (\code{too many lookups}). \\[3pt]
  \multicolumn{2}{@{}l}{\textsf{\textbf{Say hi and chat}}} \\[2pt]
  \code{send_request}\newline{\scriptsize target}\newline{\scriptsize returns ID and end time}
  & Creates a request to \code{target} that expires in five minutes. Checks, in order: signed in, with an
    active profile (\code{account restricted}); not the caller (\code{cannot request yourself}); no block either way and target active (\code{not available}).
    If the caller is already waiting on this same person, it returns that request instead, so a second tap
    sends nothing new. Then: caller not already in a chat (\code{already chatting}); target live, \code{open} and near the caller,
    who must also be live (\code{not nearby}); fewer than five requests sent in ten minutes
    (\code{too many requests}). It then marks the caller's own timed-out request as expired and refuses if
    one is still waiting (\code{already waiting}). The partial unique index enforces the same rule if two
    calls race. \\
  \code{cancel_request}\newline{\scriptsize request\_id}
  & The sender withdraws a request. Locks the caller's request if it is \code{pending} or
    \code{accepted} and sets it to \code{cancelled}; if it had just been accepted, it also returns both
    users to \code{open}. Returns silently if there is nothing left to cancel. \\
  \code{respond}\newline{\scriptsize request\_id, accept}
  & The recipient answers. Locks the request, which must be \code{pending} and addressed to the caller
    (\code{no such pending request}) and not past its end time (\code{request expired}). A decline sets
    \code{declined}. An accept requires the sender to be live (\code{not available}), withdraws the
    caller's own waiting request, refuses if the caller (\code{already chatting}) or the sender
    (\code{not available}) is already in a chat, then sets \code{accepted} and sets both users to
    \code{chatting}. \\
  \code{request_person}\newline{\scriptsize request\_id}\newline{\scriptsize returns a person card}
  & Returns the other participant's card, or nothing. Nothing if either user has blocked the other or either is not active. The
    recipient may see the sender's card at any stage; the sender may see the recipient's card only once
    the request was accepted (status \code{accepted} or \code{met}, or \code{cancelled} after an accept).
    The live note is removed once the request or chat is over. \\
  \code{end_chat}\newline{\scriptsize request\_id, outcome}
  & Ends the caller's accepted chat with outcome \code{met} or \code{cancelled} (\code{invalid outcome}
    otherwise; \code{no such chat} if it is not an accepted chat of the caller) and returns both users to
    \code{open}. \\
  \code{send_chat_message}\newline{\scriptsize request\_id, body}\newline{\scriptsize returns the message}
  & Checks: an active profile (\code{account restricted}); body of 1 to 1000 characters (\code{invalid
    message}); an accepted chat of the caller (\code{no such chat}); no block between the two (\code{not available}); fewer than 30 messages sent by
    the caller in the last minute (\code{too many messages}). Stores the message with the sender's display
    name and returns the saved row. \\
  \code{my_status}\newline{\scriptsize returns live, chat and request}
  & What the app needs after a restart: the caller's live session (note, start and end) if not expired;
    the ID of the caller's newest accepted chat in which at least one person is still live; and the
    caller's waiting request with the recipient's card, if one has not expired. Also returns the server's
    clock (\texttt{now}, in milliseconds), which the app uses to correct for a phone whose clock is wrong
    (\cref{sec:time}). \\[3pt]
  \multicolumn{2}{@{}l}{\textsf{\textbf{Safety and account}}} \\[2pt]
  \code{block_user}\newline{\scriptsize target}
  & Records the block (repeating it is harmless). Checks: signed in; not the caller
    (\code{cannot block yourself}). Cancels any accepted chat and any waiting request between the two; if
    it ended a chat, returns both users to \code{open}. A chat either user has with someone else is not
    affected. \\
  \code{blocked_list}\newline{\scriptsize returns names}
  & The caller's blocked users, newest first, as ID, first name and last initial. \\
  \code{report_user}\newline{\scriptsize target, reason, details, request\_id}
  & Reports a person. Saves a copy of their profile (name, intro, interests, live note, photo paths)
    in the report's \code{content}, and, when the request is a say-hi or chat between the two users, the
    last 20 messages the reported person sent in it; any other request is dropped. Then files the report
    through \code{private.file_report}, which every report goes through. Its checks: signed in; not the
    caller (\code{cannot report yourself}); the target exists (\code{not available}); the reason is one of
    the six values (\code{invalid reason}); the details are at most 500 characters of clean text
    (\code{invalid details}). A second report of the same thing (the same person, message or event)
    within a day returns without a new row, so a double tap files one report; otherwise fewer than ten
    reports by the caller in 24 hours (\code{too many reports}). Records whether the report counts towards
    the automatic hide (\cref{sec:blocking-reporting}). Sets the reported profile to \code{hidden}, and
    ends everything the person is in, once three different users' counted reports that are not yet
    reviewed name it; a per-person lock makes reports sent at once still count. \\
  \code{report_message}\newline{\scriptsize message\_id, reason, details}
  & Reports one group-chat message. Refused (\code{content not available}) unless the message exists
    and the caller is in its event. Files a report about its sender through \code{private.file_report},
    with \code{message_id}, \code{event_id} and a copy of the body, send time and event title. \\
  \code{report_event}\newline{\scriptsize event\_id, reason, details}
  & Reports an event. Refused (\code{content not available}) unless the caller joined it or the map would
    list it for them (not blocked, started at most 3 hours ago, within 50\,km of the position the server
    holds). Files a report about its host through \code{private.file_report}, with \code{event_id} and a
    copy of the emoji, title, description, start and creation times. \\
  \code{delete_account}
  & Refuses (\code{photos not removed}) while any file is left in the caller's photo folder. The client
    removes every file there through the Storage API before calling it. Returns the caller's chat partner,
    if any, to \code{open}, deletes any ban on the caller's email address, their upload times and their
    Auth sign-in flows and audit records, then deletes the caller's Auth user, which cascades through every
    table except \code{reports}, where the caller's id becomes null (\cref{sec:flow-delete}). \\[3pt]
  \multicolumn{2}{@{}l}{\textsf{\textbf{Map events}}} \\[2pt]
  \code{set_map_position}\newline{\scriptsize lat, lng}
  & Stores the centre of the position's cell as the caller's last known position, which the map uses for
    two hours. Checks: an active profile; latitude and longitude in range (\code{invalid location}). Applies
    the relocation rules under the same per-user lock as \code{go_green}: a move of 150\,m or less from the
    last known position keeps it and renews its time; a larger move must be within 5\,km plus 1000\,km/h
    over the time since the last known position was set, however long ago (\code{location jumped}); and a
    change counts towards the three in ten minutes (\code{too many location changes}) unless the caller's
    live position is already there. \\
  \code{events_near}\newline{\scriptsize radius\_m}\newline{\scriptsize returns event cards}
  & Returns up to 50 event cards within \code{radius_m} (default 10\,km, never more than 50\,km) of the
    position the server holds for the caller: the last position \code{set_map_position} or \code{go_green}
    stored, if set in the last two hours (\code{location needed} otherwise). Only events that started less than three hours ago or
    have not started, whose host is active and not blocked in either direction, ordered by start time.
    At most 20 calls a minute (\code{too many lookups}). \\
  \code{create_event}\newline{\scriptsize emoji, title, description, lat, lng, starts\_at\_ms}\newline{\scriptsize returns an event card}
  & Creates an event and adds the host as its first member. Checks: signed in; profile finished and adult
    (\code{profile not ready}) and active (\code{account restricted}); valid position (\code{invalid location}); a title of 1 to 60 characters on
    one line with something visible (\code{invalid title}); details of at most 200 characters of clean text
    (\code{invalid details}); start between five minutes ago and seven days ahead, checked before the
    number is converted (\code{invalid start time}); fewer than five events created by the caller in 24
    hours, counting events since deleted (\code{too many events}). Each creation is recorded in
    \code{private.event_posts}. The category must be one of the app's twelve (check constraint
    \code{events_emoji_category}). \\
  \code{join_event}\newline{\scriptsize event\_id}
  & Adds the caller as a member (repeating it is harmless). Checks: signed in; the event exists
    (\code{no such event}); no block between the caller and the host (\code{not available}); an active
    profile; a position the server holds for the caller, as for \code{events_near} (\code{location
    needed}), within 5\,km of the event (\code{too far away}). The app sends the phone's position to
    \code{set_map_position} first. \\
  \code{delete_event}\newline{\scriptsize event\_id}
  & Deletes the event; its members and group-chat messages go with it through cascading foreign keys.
    Checks: signed in; the event exists (\code{no such event}); the caller is its host (\code{not the host}). \\
  \code{leave_event}\newline{\scriptsize event\_id}
  & If the caller hosts the event, deletes it with its members and messages; otherwise removes the
    caller's membership. \\
  \code{send_event_message}\newline{\scriptsize event\_id, body}\newline{\scriptsize returns the message}
  & Checks: an active profile (\code{account restricted}); body of 1 to 1000 characters (\code{invalid
    message}); the caller is a member (\code{not a member}); fewer than 30 messages in the last minute (\code{too many messages}). Stores
    the message with the sender's display name and returns the saved row. \\
\end{longtable}
\endgroup

\Cref{tab:crud} shows which tables each function reads and writes. It is the quickest way to find every
function that can change a given table.

\begin{table}[ht]
  \centering
  \caption[Table access by function]{Table access by function: R read, I insert, U update, D delete. Columns: \code{profiles},
    \code{availability}, \code{chat_requests}, \code{messages}, \code{events}, \code{event_members},
    \code{blocks}, \code{reports}, \code{private.location_changes}.}
  \label{tab:crud}
  \tablefont
  \newcommand{\h}[1]{\rotatebox{90}{\code{#1}}}
  \setlength{\tabcolsep}{3pt}
  \begin{tabular}{@{}l*{9}{>{\centering\arraybackslash}p{8.5mm}}@{}}
    \toprule
    & \h{profiles} & \h{availability} & \h{chat_requests} & \h{messages} & \h{events} & \h{event_members} & \h{blocks} & \h{reports} & \h{location_changes} \\
    \midrule
    \code{go_green} & R & RIU & R & & & & & & RI \\
    \code{set_map_position} & R & R & & & & & & & RI \\
    \code{update_note} & R & U & & & & & & & \\
    \code{go_offline} & & D & U & & & & & & \\
    \code{nearby} & R & R & & & & & R & & \\
    \code{send_request} & R & R & RIU & & & & R & & \\
    \code{cancel_request} & & U & RU & & & & & & \\
    \code{respond} & & RU & RU & & & & & & \\
    \code{request_person} & R & R & R & & & & R & & \\
    \code{end_chat} & & U & RU & & & & & & \\
    \code{send_chat_message} & R & & R & RI & & & R & & \\
    \code{my_status} & R & R & R & & & & & & \\
    \code{block_user} & & U & U & & & & I & & \\
    \code{blocked_list} & R & & & & & & R & & \\
    \code{report_user} & RU & RUD & RU & R & & RD & & RI & \\
    \code{report_message} & RU & UD & RU & R & R & RD & & RI & \\
    \code{report_event} & RU & UD & RU & R & R & RD & R & RI & \\
    \code{delete_account} & RD\textsuperscript{*} & U & R & & & & & U & \\
    \code{events_near} & R & R & & & R & R & R & & \\
    \code{create_event} & R & & & & RI & RI & & & \\
    \code{join_event} & R & R & & & R & I & R & & \\
    \code{leave_event} & & & & & RD & D & & & \\
    \code{delete_event} & & & & D\textsuperscript{\dag} & RD & D\textsuperscript{\dag} & & & \\
    \code{send_event_message} & R & & & RI & & R & & & \\
    \code{expire_stale} (job) & & RD & RU & & & & & & D \\
    \midrule
    Client, directly & RIU & & R & R & & & D & & \\
    \bottomrule
  \end{tabular}
  \par\smallskip\raggedright\footnotesize\textsuperscript{*}\,\code{delete_account} deletes the
  \code{auth.users} row; the cascade removes the profile and every row that references it, and sets
  the user's id to null in \code{reports}. Before that it deletes the user's rows in
  \code{private.banned_identities} and \code{private.photo_uploads}, which have no foreign key.
  \textsuperscript{\dag}\,By cascade from the deleted event; \code{leave_event} by its host deletes in the same way.
  Not shown: \code{create_event} also inserts into \code{private.event_posts}; \code{set_map_position}
  and \code{go_green} write \code{private.last_positions}, which they, \code{events_near} and
  \code{join_event} read; and \code{nearby} and \code{events_near} count calls in
  \code{private.call_counts}. The minute job prunes the event posts and call counts.
\end{table}

Two details of these functions carry more weight than their size suggests. First, \code{respond},
\code{cancel_request} and \code{end_chat} lock the request row with \code{SELECT \dots FOR UPDATE} before
deciding, so a cancel and an accept that arrive together are applied one after the other, and the second
sees the first one's result. Second, \code{respond} withdraws the caller's own waiting request
\emph{before} it checks whether the caller is already chatting. If an answer to that request is being
committed at the same moment, the update waits for it, and the check that follows sees the result.

\section{Direct table access}
\label{sec:direct-access}

The client reads and writes a few tables directly through the Data API, always within the limits of
row-level security and column grants (\cref{tab:direct-access}).

\begin{table}[ht]
  \centering
  \caption{Direct table and storage access by the client.}
  \label{tab:direct-access}
  \tablefont
  \begin{tabularx}{\textwidth}{@{}L{30mm}L{40mm}Y@{}}
    \toprule
    \textsf{\textbf{Table}} & \textsf{\textbf{Operation}} & \textsf{\textbf{Used for}} \\
    \midrule
    \code{profiles} & select, insert, update own row & Loading and saving the user's own profile (\code{fetchMyProfile}, \code{saveProfile}) \\
    \code{chat_requests} & select rows the user takes part in & Missed incoming requests (\code{pendingIncoming}), the status of a request or chat (\code{requestStatus}), the answer poll in \code{waitForAnswer} \\
    \code{messages} & select & Chat and group chat history, newest 200 (\code{fetchMessages}) \\
    \code{blocks} & delete own row & Unblocking (\code{unblock}) \\
    \code{storage.objects} & upload, list, remove, sign & Photos in the user's own folder; signed URLs for photos the user may see \\
    \bottomrule
  \end{tabularx}
\end{table}

A profile row is inserted the first time it is saved and updated after that, never upserted. An upsert
would also try to set \code{id} and \code{is_adult} on conflict, and the column grants allow a client to
update only the editable columns, so the database would refuse it.

\section{Private helpers}
\label{sec:helpers}

The functions share a small set of helpers in the \code{private} schema (\cref{tab:helpers}). Those
used by a row-level security policy, a storage policy or a check constraint must be executable by
signed-in users, because the policy or constraint runs them as the user. Each answers only for the
caller, or reads no tables at all.

\begingroup\tablefont
\begin{longtable}{@{}L{42mm}L{\dimexpr\linewidth-42mm-2\tabcolsep\relax}@{}}
  \caption{Private helpers.}
  \label{tab:helpers}\\
  \toprule
  \textsf{\textbf{Helper}} & \textsf{\textbf{Purpose}} \\
  \midrule
  \endfirsthead
  \toprule
  \textsf{\textbf{Helper}} & \textsf{\textbf{Purpose}} \\
  \midrule
  \endhead
  \bottomrule
  \endfoot
  \code{blocked_between(a, b)} & True if either user has blocked the other \\
  \code{person_card(target)} & The JSON card that is all another user may see of a profile \\
  \code{near_each_other(a, b)} & True if both users are live and their cell centres are at most 150\,m apart: the same cell or a touching one \\
  \code{cell_center(lat, lng)} & The centre of the cell of about 100\,m a position falls in; the only position ever stored \\
  \code{my_position(user)} & The position the server holds for a user: the last position from the map or going live, if set in the last two hours \\
  \code{require_active(user)} & Refuses with \code{not signed in}, \code{profile not ready} or \code{account restricted} unless the user has an active adult profile \\
  \code{count_call(user, list, max, per)} & Counts a call to \code{nearby} or \code{events_near} in a fixed window and refuses past the limit (\code{too many lookups}) \\
  \code{event_card(event, viewer)} & An event's fields, host name, member count and whether the viewer has joined; the host's photo path only when the viewer has joined \\
  \code{is_chat_participant(request)} & Used by the \code{messages} select policy; executable by signed-in users \\
  \code{is_event_member(event)} & Used by the \code{messages} select policy; executable by signed-in users \\
  \code{in_own_folder(paths, owner)} & Used by the photo check constraints; reads no tables \\
  \code{has_visible_text(value)} & True if the text shows at least one character, not only spaces, line breaks or invisible characters; used by check constraints and \code{create_event}; reads no tables \\
  \code{is_clean_text(value, multiline)} & True if the text has no control characters, direction overrides or more than three stacked accents, and no line breaks unless \code{multiline}; used the same way \\
  \code{valid_interests(interests)} & True for 3 to 5 different one-line interests of 1 to 30 characters in a plain list; used by a check constraint \\
  \code{my_photo_count()} & How many files the caller has in the \code{photos} bucket; used by the storage insert policy, which cannot read \code{storage.objects} itself \\
  \code{my_uploads_today()} & How many photos the caller uploaded in the last day; used by the storage insert policy \\
  \code{record_photo_upload()} & Trigger on \code{storage.objects}: records each upload to the \code{photos} bucket in \code{private.photo_uploads} \\
  \code{can_see_photo(object name)} & Used by the storage read policy (\cref{sec:photo-privacy}); executable by signed-in users \\
  \code{has_blocked_word(input)} & True if the text contains a word from \code{private.blocked_words} (\cref{sec:word-filter}) \\
  \code{refuse_blocked_words()} & The word-filter trigger function: refuses the row when a named column has a listed word \\
  \code{end_activity(user)} & Ends a user's live session, requests, chats and event memberships; used by the automatic hide and \code{suspend_user} \\
  \code{normalised_email(email)}, \code{email_hash(email)} & An address as the ban list compares it (lowercase, no \code{+tag}, no dots for Gmail), and its SHA-256 \\
  \code{ban_email_of(user, reason)} & Puts a user's email address on \code{private.banned_identities} \\
  \code{set_reported_hash()} & Trigger on \code{reports}: stores a hash of the reported id, which outlives the account \\
  \code{expire_stale()} & The minute job (\cref{sec:housekeeping}) \\
\end{longtable}
\endgroup

The definition of \emph{near} is worth quoting, because it decides who appears in a nearby list, who can
be sent a request and whose photos a user may load:

\begin{lstlisting}[language=pgsql, caption={The proximity test in \code{private.near_each_other}.}, label={lst:near}]
select exists (
  select 1
  from public.availability av_a
  join public.availability av_b on av_b.user_id = b
  where av_a.user_id = a
    and av_a.expires_at > now()
    and av_b.expires_at > now()
    and extensions.st_dwithin(av_a.location, av_b.location, 150)
);
\end{lstlisting}

Both stored positions are cell centres, so the test asks whether the two users are in the same cell of
about 100 metres or a touching one. The distance is fixed and the same for everyone; nothing the caller
sends changes it (\cref{sec:location-privacy}). Distances are computed on the geography type, in metres
on the Earth's surface.

\section{Word filter}
\label{sec:word-filter}

The app stores require a filter for objectionable content in an app where people post text and meet.
Yap's filter is a short list of words in \code{private.blocked_words}, checked by triggers before a row is
saved (\cref{tab:word-filter}). A row that contains a listed word is refused with \code{objectionable
text}, which the app shows as \ui{That has words Yap doesn't allow (see the community guidelines). Change
them and try again.} Because the check runs on the tables, it covers the profile, which the client writes
directly, as well as the database functions that save a note, an event or a message.

\begin{table}[ht]
  \centering
  \caption{Text checked by the word filter.}
  \label{tab:word-filter}
  \tablefont
  \begin{tabularx}{\textwidth}{@{}L{34mm}L{44mm}Y@{}}
    \toprule
    \textsf{\textbf{Trigger}} & \textsf{\textbf{Columns}} & \textsf{\textbf{Written by}} \\
    \midrule
    \code{profiles_word_filter} & \code{intro}, \code{interests} & The client's own profile saves \\
    \code{availability_word_filter} & \code{note} & \code{go_green}, \code{update_note} \\
    \code{events_word_filter} & \code{title}, \code{description} & \code{create_event} \\
    \code{messages_word_filter} & \code{body} & \code{send_chat_message}, \code{send_event_message} \\
    \bottomrule
  \end{tabularx}
\end{table}

\code{private.has_blocked_word} lowercases the text, splits it into words at every character that is not a
letter or digit, and compares each word with the list exactly. It repeats the comparison once with common
look-alike characters swapped back (\code{0} to \code{o}, \code{1} to \code{i}, \code{3} to \code{e},
\code{4} to \code{a}, \code{5} to \code{s}, \code{7} to \code{t}, \code{@} to \code{a}, \texttt{\$} to
\code{s}), so \code{h4te} is checked as \code{hate}. Only whole words match, so a longer word that merely
contains a listed one is never refused. Interests, a text array, are checked as their JSON text, which
splits the same way.

Names are not checked, because a real name can match a listed word, and report details are not checked,
so people can quote what was said to them. The list is deliberately narrow: it leaves out general
swearing, sexual words, words people use to describe harm or ask for help, and words that are also names
or ordinary words. It is a first line only; reports and blocks handle everything else. Rows saved before
the filter existed are not checked again; only new text and edits are.

\section{Moderation function}
\label{sec:moderation}

\code{private.suspend_user(target, note)} acts on a report after the owner has reviewed it. It is run from
the Supabase SQL editor and cannot be called from the app: execution is revoked from every API role. It
refuses an unknown profile (\code{no such profile}), then:

\begin{enumerate}
  \item sets the profile's status to \code{suspended}, so the person cannot go live, say hi, send
    messages or create or join events (\code{account restricted}), and nobody sees their profile, photos
    or events;
  \item through \code{private.end_activity}, ends every chat they are in, returning the other person to
    \code{open} if still live, cancels their waiting requests, and deletes their live session and event
    memberships, so they cannot post in group chats either;
  \item puts their email address on the ban list until the suspension is lifted or the account is
    deleted, so signing up again with it is refused (\cref{sec:signup-checks}).
  \item marks every open report about them as reviewed, with \code{note} (default \code{suspended}) as
    the action.
\end{enumerate}

\code{private.restore_user(target)} lifts a suspension or an automatic hide: it sets the status back to
\code{active}, makes the person a member again of the events they still host (\code{end_activity} took
them out), takes the email address off the ban list, and marks every open report about the person
as reviewed (action \code{restored}), so reports already reviewed cannot hide them again. \code{docs/legal/moderation.md} describes
the owner's daily review: the queries that list open reports and what the reported person posted, and
the statements for each outcome, from marking reports reviewed and showing a hidden profile again, through
removing one message or event, to suspending or deleting the account.

\section{Scheduled housekeeping}
\label{sec:housekeeping}

The \code{pg_cron} job \code{yap-expire-stale} runs \code{private.expire_stale()} every minute. It makes
time-based rules hold even when no app is open:

\begin{enumerate}
  \item It deletes every live session whose end time has passed.
  \item It marks as \code{expired} every \code{pending} request whose five minutes are up or whose sender
    is no longer live.
  \item It marks as \code{cancelled} every \code{accepted} chat in which neither person is live any more.
  \item It deletes entries older than ten minutes from \code{private.location_changes}.
  \item It deletes entries older than a day from \code{private.event_posts},
    \code{private.signup_ips} and \code{private.photo_uploads}, and call counts from windows more than an
    hour old. Last known positions are kept.
\end{enumerate}

Because the job runs once a minute, a row can outlive its end time by up to a minute. The functions do
not rely on the job for correctness: they compare end times with \code{now()} themselves, and
\code{send_request} expires the caller's own timed-out request before it counts waiting requests.

\section{Realtime publication}
\label{sec:realtime-publication}

The tables \code{chat_requests} and \code{messages} belong to the \code{supabase_realtime} publication,
and both have \code{REPLICA IDENTITY FULL}, so Realtime can evaluate the select policies on updates as
well as inserts. A client receives an insert or update on these tables only if it could select the row
itself: requests it sends or receives, and messages of chats it takes part in or events it has joined,
excluding messages from users it has blocked. \Cref{sec:realtime} describes the channels the client opens.
